% !TeX spellcheck = it_IT
% !TeX root = UsoMail.tex
\chapter{Microsoft 365}
%\printpartialtoc
\section{Introduzione}
Quello che segue non appunti disordinati di ciò che è stato fatto in questi giorni, più un diario (incompleto) che una guida.

La Scuola ha ottenuto 2000 licenze A1 di Microsoft365\index{Microsoft365} ex Office365\index{Office365} per i docenti e altrettante per gli studenti.
\section{Domini}
Con 365 la scuola ha due domini:
\begin{itemize}
	\item IISCavourMarconiPascal981.onmicrosoft.com, ‎Gestito in ‎Microsoft 365‎‎ - ‎Dominio predefinito, Dominio di fallback‎. ‎Microsoft‎ fornisce un dominio ‎onmicrosoft.com‎ di fallback nel caso tu non sia proprietario di un dominio o non voglia connetterlo a ‎Microsoft 365‎. Il dominio di fallback viene utilizzato per impostazione predefinita in:
	\begin{itemize}
		\item Nomi utente e indirizzi di posta elettronica
		\item ‎Microsoft 365‎ Team e gruppi alias di posta elettronica
		\item Spostamenti automatici delle dipendenze di dominio
	\end{itemize}
	Per modificare il dominio di fallback, aggiungere un altro dominio ‎onmicrosoft.com‎ con un nome personalizzato, quindi impostarlo come fallback\cite{Microsoft2023}
	\item iisperugia.net ‎Gestito in ‎enomCentral‎‎ - ‎Dominio federato‎\cite[postnote]{Google2023h}
\end{itemize} 

\section{Federation SSO}
L'introduzione di 365 porta a problemi di gestione e creazione ex novo di utenti e password\cite{Google2023i}.
\section{Gruppi}
Creato il gruppo \textbf{docenti}
\section{Licenze}
Le licenze si assegnano tramite Fatturazione\textrightarrow Licenze \textrightarrow Assegna le licenze 

Si possono inserire venti utenti per volta. 